\section{闭球的紧性}
\label{sec:lfunc-ball-compactness}

有限维空间中，有界序列总有收敛子列，因为闭球是紧的。无限维空间中闭球在范数下不紧，这是第~\ref{chap:functions-as-vectors}章的结论。换成上一节较弱的拓扑，紧性可以恢复，但要分清三件事：闭球在哪个空间里，用的是哪个拓扑，以及紧性能否转化为“有界序列有收敛子列”。本节先证明对偶空间的闭球在弱星拓扑下紧，再借助反射性把它转回原空间闭球的弱紧性，最后讨论收敛子列的提取。

记$B_X=\{x\in X:\lVert x\rVert\leq1\}$，$B_{X^*}=\{\ell\in X^*:\lVert\ell\rVert\leq1\}$，分别是原空间与对偶空间的闭单位球。

\subsection{对偶闭球的弱星紧性}
\label{sec:lfunc-banach-alaoglu}

弱星拓扑只要求逐个向量上的读数接近，测试条件很少，开集也就很少；开集越少，开覆盖越难找到无法取有限子覆盖的反例，紧性越容易成立。下面的定理说，对偶空间的闭单位球在这个拓扑下总是紧的。

\begin{theorem}[Banach--Alaoglu定理]
\label{thm:lfunc-banach-alaoglu}
设$X$是赋范空间。则$B_{X^*}$在弱星拓扑$\sigma(X^*,X)$下紧。
\end{theorem}

这一结论称为\term{Banach--Alaoglu定理}（Banach--Alaoglu theorem）。一般情形的证明把每个$\ell\in B_{X^*}$看成一族数$(\ell(x))_{x\in X}$，其中$\lvert\ell(x)\rvert\leq\lVert x\rVert$，从而把$B_{X^*}$放进无穷多个闭圆盘的乘积中；弱星拓扑恰好是乘积拓扑在其上的限制，线性条件$\ell(\alpha x+\beta y)=\alpha\ell(x)+\beta\ell(y)$又使$B_{X^*}$成为乘积中的闭集。剩下的一步需要“任意多个紧空间的乘积仍紧”的Tychonoff定理，它超出了第~\ref{chap:topological-spaces}章所讨论的有限乘积，这里不给出一般情形的证明。$X$可分时，第~\ref{sec:lfunc-subsequences}小节将用度量和对角线论证给出完整证明。

\paragraph{推广到任意闭球}
对$\ell_0\in X^*$和$r>0$，映射$T(\ell)=\ell_0+r\ell$把$B_{X^*}$映成闭球$\{\ell:\lVert\ell-\ell_0\rVert\leq r\}$。它把弱星邻域$U^*(\ell;x_1,\ldots,x_m;\varepsilon)$恰好映成$U^*(T\ell;x_1,\ldots,x_m;r\varepsilon)$，逆映射同理，所以$T$是弱星拓扑下的同胚。紧集在连续映射下的像是紧的，于是$X^*$中任意有限半径的闭球都是弱星紧的。

\paragraph{定理断言了什么}
定理不要求$X$完备，也不要求任何其他条件。它断言的是对偶空间$X^*$的闭球在$\sigma(X^*,X)$下紧，既不是原空间$X$的闭球在弱拓扑下紧，也不是$B_{X^*}$在范数下紧。后者一般不成立：取$X=\ell^1$，$X^*=\ell^\infty$，$\ell^\infty$中的单位向量$\bm e_n$两两距离为$1$，没有范数收敛的子列，$B_{\ell^\infty}$在范数下不紧。

\subsection{反射性}
\label{sec:lfunc-reflexivity}

要把对偶空间上的结论带回原空间，需要一座桥。每个$x\in X$都给出$X^*$上的一个读数：用$x$去读一个泛函$\ell$，得到$\ell(x)$。

\begin{definition}[自然嵌入与反射空间]
\label{def:lfunc-reflexive}
对赋范空间$X$，令$J:X\to X^{**}$，
\[
J(x)(\ell)=\ell(x)\qquad(x\in X,\ \ell\in X^*),
\]
称为\term{自然嵌入}（canonical embedding）。若$J$是满射，称$X$为\term{反射空间}（reflexive space）。
\end{definition}

$J(x)$对$\ell$线性，且$\lvert J(x)(\ell)\rvert\leq\lVert\ell\rVert\lVert x\rVert$，所以$J(x)\in X^{**}$；$J$本身对$x$线性。由范数的对偶表示~\eqref{eq:lfunc-norm-dual}，
\[
\lVert J(x)\rVert=\sup_{\lVert\ell\rVert\leq1}\lvert\ell(x)\rvert=\lVert x\rVert,
\]
所以$J$是等距的，特别是单射。反射空间中$J$是等距同构，$X$与完备的$X^{**}$（定理~\ref{thm:lfunc-dual-complete}）等距对应，所以反射空间自动是Banach空间。定义要求的是这个特定的$J$为满射；仅仅存在某个从$X$到$X^{**}$的等距同构并不够，这里不展开相应的反例。

\paragraph{Hilbert空间是反射的}
记$R:H\to H^*$，$R(h)=\langle h,\cdot\rangle$。由Riesz表示定理，$R$是保范数的双射，并且共轭线性。给定$\Phi\in H^{**}$，映射$h\mapsto\overline{\Phi(Rh)}$对$h$线性（两次共轭抵消）且有界，再用一次Riesz表示定理，存在$x\in H$使$\overline{\Phi(Rh)}=\langle x,h\rangle$对一切$h$成立。于是
\[
\Phi(Rh)=\overline{\langle x,h\rangle}=\langle h,x\rangle=(Rh)(x)=J(x)(Rh).
\]
$R$是满射，所以$\Phi=J(x)$。

\paragraph{\texorpdfstring{$L^p$}{Lp}空间（\texorpdfstring{$1<p<\infty$}{1<p<infinity}）是反射的}
设测度$\sigma$有限，$q$为共轭指数，则$1<q<\infty$。由定理~\ref{thm:lfunc-lp-dual}，$(L^p)^*$的每个元素都是某个$L_g$，$g\in L^q$。给定$\Phi\in(L^p)^{**}$，映射$g\mapsto\Phi(L_g)$是$L^q$上的有界线性泛函，再对$L^q$用定理~\ref{thm:lfunc-lp-dual}，存在$f\in L^p$使$\Phi(L_g)=\int gf\,d\mu=L_g(f)=J(f)(L_g)$对一切$g$成立，所以$\Phi=J(f)$。两次表示都需要指数严格小于$\infty$，这就是端点被排除的原因。

\paragraph{完备而不反射}
$\ell^1$是Banach空间，却不反射。把$(\ell^1)^*$等同于$\ell^\infty$，则$J(x)(y)=\sum_ky_kx_k$。第~\ref{sec:lfunc-hahn-banach}小节的$\Lambda\in(\ell^\infty)^*=(\ell^1)^{**}$若等于某个$J(x)$，则$x_k=J(x)(\bm e_k)=\Lambda(\bm e_k)=0$对一切$k$成立，$x=0$，与$\Lambda(1,1,\ldots)=1$矛盾。所以完备性只说明Cauchy列有极限，反射性则要求二次对偶中没有多余的元素，后者严格更强。第~\ref{sec:lfunc-subsequences}小节还会看到$L^1(0,1)$也不反射。

\subsection{反射空间中的弱紧性}
\label{sec:lfunc-reflexive-weak-compact}

设$X$反射。$J$是等距的满射，所以它把$B_X$恰好映成$B_{X^{**}}$：$\lVert x\rVert\leq1$当且仅当$\lVert J(x)\rVert\leq1$，而$B_{X^{**}}$中每个元素都是某个$J(x)$。

\paragraph{两种拓扑的对应}
$X$上的弱拓扑$\sigma(X,X^*)$用有限个$\ell_i\in X^*$测试，$X^{**}=(X^*)^*$上的弱星拓扑$\sigma(X^{**},X^*)$也用有限个$\ell_i\in X^*$测试。由$J(x)(\ell_i)=\ell_i(x)$，
\[
J\bigl(U(x_0;\ell_1,\ldots,\ell_m;\varepsilon)\bigr)=U^*\bigl(J(x_0);\ell_1,\ldots,\ell_m;\varepsilon\bigr),
\]
右端用到了$J$为满射，否则只能得到右端与$J(X)$的交。两种拓扑的基本邻域一一对应，所以$J$是从$(X,\sigma(X,X^*))$到$(X^{**},\sigma(X^{**},X^*))$的同胚。

\begin{theorem}[反射空间中闭球的弱紧性]
\label{thm:lfunc-reflexive-weak-compact}
设$X$是反射空间。则$B_X$在弱拓扑$\sigma(X,X^*)$下紧；$X$中任意有限半径的闭球也是弱紧的。
\end{theorem}

\begin{proof}
对赋范空间$X^*$应用定理~\ref{thm:lfunc-banach-alaoglu}，$B_{X^{**}}$在$\sigma(X^{**},X^*)$下紧。由上面的对应，$B_X=J^{-1}(B_{X^{**}})$是这个紧集在同胚$J^{-1}$下的像，所以在弱拓扑下紧。平移与缩放$x\mapsto x_0+rx$在弱拓扑下是同胚，理由与第~\ref{sec:lfunc-banach-alaoglu}小节相同，所以任意闭球弱紧。
\end{proof}

由此，Hilbert空间和$1<p<\infty$时的$L^p$空间中，闭球都是弱紧的。这一结论不保证范数紧性：$\ell^2$的闭单位球弱紧，但$\bm e_n$没有范数收敛的子列。弱紧性换来的是更弱的收敛方式，不是无限维困难的消失。

\subsection{收敛子列的提取}
\label{sec:lfunc-subsequences}

紧性说的是开覆盖有有限子覆盖，我们真正想用的却常常是“有界序列有收敛子列”。在度量空间中两者等价，这是第~\ref{chap:metric-spaces}章的结论；但第~\ref{chap:topological-spaces}章已经说明，一般拓扑空间中紧性不能直接推出序列紧性。弱拓扑和弱星拓扑在无限维时一般不能由度量给出，所以上面的紧性定理还不能直接用来提取子列。

\paragraph{可分前对偶下弱星拓扑可度量}
设$X$可分，$\{w_k\}$是$X$的可数稠密子集。令$z_k=w_k/\max(1,\lVert w_k\rVert)$，则$\{z_k\}\subset B_X$，并且在$B_X$中稠密：若$\lVert x\rVert\leq1$且$\lVert x-w_k\rVert<\varepsilon$，则$\lVert w_k\rVert<1+\varepsilon$，从而$\lVert z_k-w_k\rVert<\varepsilon$，$\lVert x-z_k\rVert<2\varepsilon$。在$B_{X^*}$上令
\[
d(\ell,\ell')=\sum_{k=1}^{\infty}2^{-k}\lvert\ell(z_k)-\ell'(z_k)\rvert .
\]
每项不超过$2\cdot2^{-k}$，级数收敛。对称性与三角不等式逐项成立；若$d(\ell,\ell')=0$，则$\ell$与$\ell'$在$B_X$的稠密子集上相等，由连续性在$B_X$上相等，由线性在$X$上相等。所以$d$是$B_{X^*}$上的度量。

\begin{lemma}[球上弱星拓扑的度量]
\label{lem:lfunc-weak-star-metric}
设$X$可分。则在$B_{X^*}$上，度量$d$诱导的拓扑与弱星拓扑$\sigma(X^*,X)$限制到$B_{X^*}$上的拓扑相同。
\end{lemma}

\begin{proof}
只需说明每一点处，两种拓扑的邻域互相包含。

\emph{度量球包含弱星邻域。}给定$\ell_0\in B_{X^*}$和$r>0$，取$K$使$2^{1-K}<r/2$。若$\ell\in B_{X^*}$且$\lvert\ell(z_k)-\ell_0(z_k)\rvert<r/2$对$k\leq K$成立，则$d(\ell,\ell_0)<\sum_{k\leq K}2^{-k}r/2+\sum_{k>K}2\cdot2^{-k}<r/2+r/2=r$。

\emph{弱星邻域包含度量球。}给定$U^*(\ell_0;x_1,\ldots,x_m;\varepsilon)$，不妨设各$x_i\neq0$，令$c=\max_i\lVert x_i\rVert$，$u_i=x_i/c\in B_X$，只需使$\lvert\ell(u_i)-\ell_0(u_i)\rvert<\varepsilon/c$。对每个$i$取$k_i$使$\lVert u_i-z_{k_i}\rVert<\varepsilon/(4c)$。对$\ell\in B_{X^*}$，
\begin{equation*}
\begin{aligned}
\lvert\ell(u_i)-\ell_0(u_i)\rvert&\leq\lvert\ell(u_i-z_{k_i})\rvert+\lvert\ell(z_{k_i})-\ell_0(z_{k_i})\rvert+\lvert\ell_0(z_{k_i}-u_i)\rvert\\
&<\frac{\varepsilon}{2c}+\lvert\ell(z_{k_i})-\ell_0(z_{k_i})\rvert .
\end{aligned}
\end{equation*}
令$K=\max_ik_i$，$\delta=2^{-K}\varepsilon/(2c)$。若$d(\ell,\ell_0)<\delta$，则$\lvert\ell(z_{k_i})-\ell_0(z_{k_i})\rvert\leq2^{k_i}d(\ell,\ell_0)<\varepsilon/(2c)$，于是$\lvert\ell(u_i)-\ell_0(u_i)\rvert<\varepsilon/c$对每个$i$成立。
\end{proof}

统一的范数界在证明中起了关键作用：$\ell$与$\ell_0$都在单位球中，才能用有限个$z_k$上的接近控制任意$u_i$上的接近。在整个$X^*$上，这个度量一般不能给出弱星拓扑。

\begin{theorem}[可分前对偶下的弱星子列]
\label{thm:lfunc-weak-star-subsequence}
设$X$是可分赋范空间。则$X^*$中每个有界序列都有弱星收敛的子列。
\end{theorem}

\begin{proof}
设$\lVert\ell_n\rVert\leq C$，$\{z_k\}$如上。对每个$k$，$(\ell_n(z_k))_n$是有界数列。先抽子列使在$z_1$上的值收敛，再从中抽子列使在$z_2$上的值收敛，如此继续；取第$j$个子列的第$j$项组成对角子列$(\ell_{n_j})$，它从某项起是每个子列的子列，所以$\ell_{n_j}(z_k)$对每个$k$都收敛。

对$x\in B_X$及$\varepsilon>0$，取$z_k$使$\lVert x-z_k\rVert<\varepsilon$，则
\[
\lvert\ell_{n_i}(x)-\ell_{n_j}(x)\rvert\leq2C\varepsilon+\lvert\ell_{n_i}(z_k)-\ell_{n_j}(z_k)\rvert,
\]
最后一项在$i,j$充分大时小于$\varepsilon$。所以$(\ell_{n_j}(x))$是Cauchy数列，由线性对一切$x\in X$都收敛，极限记为$\ell(x)$。$\ell$是线性的，且$\lvert\ell(x)\rvert\leq C\lVert x\rVert$，故$\ell\in X^*$，$\ell_{n_j}\overset{*}{\rightharpoonup}\ell$。
\end{proof}

取$C=1$，极限仍在$B_{X^*}$中，所以$B_{X^*}$在度量$d$下序列紧；由第~\ref{chap:metric-spaces}章，序列紧的度量空间是紧的，再由引理~\ref{lem:lfunc-weak-star-metric}，这就是$B_{X^*}$的弱星紧性。可分时定理~\ref{thm:lfunc-banach-alaoglu}由此得到完整证明。

\paragraph{可分性不能省}
若前对偶不可分，序列结论可以失效，尽管Banach--Alaoglu定理仍保证紧性。$\ell^\infty$不可分；在$(\ell^\infty)^*$中考虑$\ell_n(x)=x_n$，$\lVert\ell_n\rVert=1$。对任意子列$(\ell_{n_j})$，取$x\in\ell^\infty$使$x_{n_j}=(-1)^j$、其余坐标为零，则$\ell_{n_j}(x)=(-1)^j$不收敛。所以这个有界序列没有弱星收敛的子列。紧与序列紧在这里确实分开了。

\paragraph{反射空间中的弱子列}
对弱拓扑有一个不需要可分性的结论。Eberlein--\v{S}mulian定理断言：在Banach空间中，一个子集的弱闭包弱紧，当且仅当该子集中的每个序列都有在$X$中弱收敛的子列。它的证明较长，这里不给出证明。与定理~\ref{thm:lfunc-reflexive-weak-compact}合起来得到：

\begin{theorem}[反射空间中的弱收敛子列]
\label{thm:lfunc-reflexive-subsequence}
设$X$是反射空间。则$X$中每个范数有界的序列都有弱收敛的子列。
\end{theorem}

\begin{proof}
有界序列含于某个闭球$B=\{x:\lVert x-x_0\rVert\leq r\}$。由式~\eqref{eq:lfunc-norm-dual}，$B=\bigcap_{\lVert\ell\rVert\leq1}\{x:\lvert\ell(x)-\ell(x_0)\rvert\leq r\}$，每个集合都是弱闭的，所以$B$弱闭，它的弱闭包就是自身；由定理~\ref{thm:lfunc-reflexive-weak-compact}，$B$弱紧。由Eberlein--\v{S}mulian定理，$B$中每个序列都有弱收敛的子列。
\end{proof}

于是在Hilbert空间和$1<p<\infty$的$L^p$空间中，有界序列总有弱收敛子列，不必另加可分性条件。$\bm e_n$本身就弱收敛到$0$；有界序列$(-1)^n\bm e_1$不收敛，但它的偶数项子列收敛。

\paragraph{非反射空间中可以失败}
在$L^1(0,1)$中令$f_n=n\mathbf{1}_{(0,1/n)}$，$\lVert f_n\rVert_1=1$。设某个子列$f_{n_j}\rightharpoonup f\in L^1$。对$0<a<1$，$g=\mathbf{1}_{(a,1)}\in L^\infty$给出$L^1$上的连续读数，而$n_j>1/a$时$\int_a^1f_{n_j}=0$，所以$\int_a^1f=0$。令$a\to0$，由控制收敛得$\int_0^1f=0$。但读数$g\equiv1$给出$\int_0^1f_{n_j}=1$，又得$\int_0^1f=1$，矛盾。所以这个有界序列没有弱收敛的子列。

质量始终为$1$，却向左端点越挤越窄；远离原点的读数看到零，整体积分看到一，没有$L^1$函数能同时满足。由定理~\ref{thm:lfunc-reflexive-subsequence}的逆否命题，$L^1(0,1)$不是反射空间。若把$f_n\,dt$看成$C([0,1])$上的读数$u\mapsto\int_0^1uf_n$，它们在逐个连续函数上收敛到$u(0)$：在更大的对偶空间中，极限以点质量的形式出现了。
